\scanpage{65}
que \inserted{$u$ soit une isogénie et que} $N$ soit premier à $l$ (où $l\ne\operatorname{car}k$) il f. et s. que le composé
\[
H^1(Y,\mathbb Z_l)\unsure{\mathord{-}|\longrightarrow} H^{2n-1}(X,\mathbb Z_l)(n-1)\mathord{-}
\frac{H^{2n-1}(X,\mathbb Z_l)(n-1)}{\text{torsion}}
\]
soit \deleted{surjectif} bijectif; \deleted{et}\inserted{Pour que de plus $H^{2n-1}(X,\mathbb Z_l)$ soit sans torsion,} il suffit pour ceci que $H^1(Y,\mathbb Z_l)\to H^{2n-1}(X,\mathbb Z_l)(n-1)$ soit bijectif.

% slanted
\margin{\uline{Th} ? C'est un cas particulier important du ``Lefschetz \unsure{vache|vachement}'' \unsure{donnant|dérivant} Lefschetz des surfaces\ldots.}
\uline{Corollaire} 2.6. Soit $X$ projective lisse conn. de \unsure{dim} $n\geq2$, et $\xi_1,\ldots,\xi_{n-1}\in\inserted{\operatorname{Pic}(X)}$ \deleted{\unsure{classes des cat. ly lin} \unsure{polarisat|provenant}} \deleted{\unsure{une telle}} des faisceaux amples, \inserted{$\zeta_1,\ldots,\zeta_{n-1}\in A^1(X)$ leurs classes.} Alors pour tout entier $l$, l'hom de cup-produit par
\[
\deleted{\unsure{\xi}}\eta=\zeta_1\zeta_2\ldots\zeta_{n-1}:
\]
\[
H^1(X,\deleted{\unsure{Z}}Q_l)\longrightarrow H^{2n-1}(X,Q_l)(n-1)
\]
est un \uline{isomorphisme}. L'isom inverse $L$ est définissable par un élément de $\deleted{\unsure{Z^1_{\mathrm{rat}}}}A^1(X\times X)\otimes Q$, \deleted{i.e. de} ou \unsure{même|encore} de $\operatorname{Cl}(X,X)\otimes Q$, et \deleted{cet élément} \inserted{\unsure{ce|cet}} dernier élément est unique et ne dépend pas de $l$.

Il faut prouver que l'hom de VA
\[
u:B_X\longrightarrow A_X
\]
% archive sheet number: 64

